% Part: normal-modal-logic
% Chapter: axioms-systems
% Section: introduction

\documentclass[../../../include/open-logic-section]{subfiles}

\begin{document}

\olfileid{nml}{prf}{int}

\olsection{مقدمة}

لدينا دلالة للغة الجهية الأساسية تُصاغ بواسطة النماذج الجهية، ولدينا
مفهوم كون !!a{formula} صالحة---أي صادقة في جميع العوالم في جميع
النماذج---أو صالحة بالنسبة إلى فئة من النماذج أو الأطر---أي صادقة في
جميع العوالم في جميع نماذج الفئة، أو في جميع النماذج المبنية على
الإطار. ويربط المنطق عادة هذه التوصيفات الدلالية للصلاحية بمفهوم
برهاني هو !!{derivability}. والغاية هي تعريف مفهوم !!{derivability}
في نسق ما بحيث تكون !!a{formula} !!{derivable} إذا وفقط إذا كانت صالحة.

أبسط نظم !!{derivation} وأقدمها تاريخيا هي ما يسمى نظم !!{derivation}
من نمط هلبرت، أو النظم البديهية. ومن السهل نسبيا بناء نظم !!{derivation}
من نمط هلبرت لكثير من المنطقات الجهية: فهي بسيطة بوصفها موضوعات
للدراسة الميتانظرية (مثلا لإثبات السلامة والتمام). غير أن استعمالها
لإثبات !!{formula}s أصعب بكثير من استعمال نظم الاستنتاج الطبيعي مثلا.

في نظم !!{derivation} من نمط هلبرت، يكون !!a{derivation} لـ!!a{formula}
متتالية من !!{formula}s تبدأ من بديهيات معينة وتصل، بواسطة عدد قليل
من قواعد الاستدلال، إلى !!{formula} المقصودة. ولأننا نريد أن يطابق
نسق !!{derivation} الدلالة، فعلينا أن نضمن صدق مجموعة الصيغ
!!{derivable} في جميع النماذج (أو في جميع النماذج التي تصدق فيها كل
البديهيات). سنعزل أولا بعض خصائص المنطقات الجهية اللازمة لذلك، وهي
المنطقات الجهية «النظامية». ولا يلزم افتراض سوى قاعدتي استدلال في
المنطقات الجهية النظامية: القياس الاستثنائي والإلزام. ونتخذ بديهيات
جميع أمثلة إحلال القضايا الصادقة صدقا تحليليا، ومعها، بحسب المنطق
الجهوي الذي نعالجه، عددا من البديهيات الجهية. وحتى إذا لم نكن نهتم إلا
بفئة جميع النماذج، فعلينا أيضا عد جميع أمثلة إحلال~\Ax{K}
\iftag{notprvDiamond}{}{و~$\Ax{Dual}$} بديهيات. وهذا وحده يولد المنطق
الجهوي النظامي الأدنى~$\Log K$.

\begin{defn}
قاعدة \emph{القياس الاستثنائي} هي مخطط الاستدلال
\begin{prooftree}
\AxiomC{$!A$}
\AxiomC{$!A \lif !B$}
\RightLabel{\MP}
\BinaryInfC{$!B$}
\end{prooftree}
نقول إن !!a{formula}~$!B$ \emph{تلزم عن}~!!{formula}s $!A$، $!C$
بالقياس الاستثنائي إذا وفقط إذا كان $!C \ident !A \lif !B$.
\end{defn}

\begin{defn}
قاعدة \emph{الإلزام} هي مخطط الاستدلال
\begin{prooftree}
\AxiomC{$!A$}
\RightLabel{\Nec}
\UnaryInfC{$\Box !A$}
\end{prooftree}
نقول إن !!{formula}~$!B$ تلزم عن !!{formula}s~$!A$ بالإلزام إذا وفقط
إذا كان $!B \ident \Box !A$.
\end{defn}

\begin{defn}
\emph{!!{derivation}} من مجموعة البديهيات~$\Sigma$ هو متتالية من
!!{formula}s $!B_1$، $!B_2$، \dots، $!B_n$، بحيث يكون كل $!B_i$ إما
\begin{enumerate}
\item مثال إحلال لقضية صادقة صدقا تحليليا، أو
\item مثال إحلال لـ!!a{formula} في~$\Sigma$، أو
\item لازما بالقياس الاستثنائي عن !!{formula}s اثنتين $!B_j$، $!B_k$
  حيث $j$، $k < i$، أو
\item لازما بالإلزام عن !!a{formula}~$!B_j$ حيث $j < i$.
\end{enumerate}
إذا وجد !!a{derivation} كهذا وكان $!B_n \ident !A$، قلنا إن
$!A$ \emph{!!{derivable} من $\Sigma$}، ورمزنا إلى ذلك بـ$\Sigma \Proves !A$.
\end{defn}

سيتبين بهذا التعريف أن مجموعة الصيغ !!{derivable} تؤلف منطقا جهيا
نظاميا، وأن كل عنصر !!{derivable} من نوع !!{formula} صادق في كل نموذج تصدق فيه
جميع البديهيات. وتسمى خاصية !!{derivation}s هذه \emph{السلامة}. أما
العكس، أي \emph{التمام}، فإثباته أصعب.

\end{document}
